\section{Introduction}

\textit{\textbf{High-Order Hybrid-Structured}} (HOHS) data, characterized by globally unstructured meshes with locally high-order, quasi-structured elements, have emerged as an important class of scientific simulation data. They combine geometric flexibility with the accuracy of high-order discretizations and arise across a broad range of scientific and engineering domains, including computational fluid dynamics, computational electromagnetics, seismology, thermosphere--ionosphere modeling, and extended magnetohydrodynamics. As computing capabilities continue to scale, HOHS datasets are rapidly growing along multiple dimensions, including the number of time steps, physical variables, spatial resolution, and simulated domain size, with individual HOHS data files already reaching petabyte scale. This rapid growth makes efficient storage and transfer of HOHS data a critical data-management challenge.

Error-bounded lossy compression~\cite{liang2018error,liu2024high,zhao2021optimizing,liang2022sz3,lindstrom2014fixed,li2023lossy,di2025survey,huang2023cuszp,gong2023mgard} has proven effective in addressing similar challenges in scientific data management by substantially reducing storage and transfer costs, making it a practical alternative to lossless compression for floating-point simulation data. Meanwhile, compared with conventional lossy codecs for floating-point simulation data, it enforces user-specified error bounds to provide predictable control over reconstruction accuracy.
Most existing error-bounded compressors are optimized for dense, structured tensor data defined on regular grids. These methods rely on spatially regular data layouts to exploit predictable neighborhood relationships, enabling high compression ratios with low metadata overhead. Recent studies have also explored error-bounded compression for unstructured scientific data, where data points are irregularly distributed and connected through mesh or graph structures~\cite{wu2025enabling,gong2024general,ren2024prediction}. These methods typically leverage spatial proximity or mesh connectivity to construct local predictors for irregular data layouts.

Despite these advances, neither structured-grid nor unstructured-data compressors fully match the hybrid organization of HOHS data. Structured-grid compressors assume global regularity, which is absent because HOHS elements are connected in an unstructured manner. Unstructured compressors, on the other hand, are designed to handle irregular connectivity, but they typically exploit only limited local neighborhood information and do not capture the quasi-structured patterns within each high-order element. As a result, directly applying existing error-bounded compressors to HOHS data often leads to suboptimal compression ratios or throughput. This mismatch leaves HOHS data without an effective error-bounded reduction mechanism for practical storage and data movement.

In this work, we bridge this gap by developing two complementary error-bounded compressors for HOHS data management. SZ3-HOHS targets both high storage reduction and high throughput, whereas MRCZ targets maximum storage reduction.
Specifically, SZ3-HOHS adapts the interpolation-based prediction framework of SZ3~\cite{liang2022sz3,zhao2021optimizing} to HOHS datasets by dynamically adjusting interpolation coefficients to account for the non-uniform spatial spacing between data points. In comparison, MRCZ leverages model-based representations to capture locally structured patterns within each data element. By decomposing simulation fields into integer-valued coefficient tensors and residual error tensors, MRCZ transforms the original data into a representation that is significantly more amenable to compression, thereby enabling much higher compression ratios.
Together, these two solutions provide different trade-offs among compression throughput, storage reduction, and end-to-end data movement.

We evaluate SZ3-HOHS and MRCZ primarily from two data-management perspectives: storage reduction and end-to-end data movement through Globus~\cite{chard2017globus}. We further study strong-scaling throughput to assess parallel efficiency on multi-CPU platforms, and visualize the reconstruction errors introduced by data reduction to examine their impact on scientific data quality. Our evaluation uses a collection of HOHS datasets~\cite{yuan2024coupled,fischer2022nekrs}, including simulations of pebble-bed reactors, crossflow jets, and laminar flows. We compare the proposed methods against state-of-the-art error-bounded lossy compressors, including ZFP~\cite{lindstrom2014fixed}, SPERR~\cite{li2023lossy}, SZ2~\cite{tao2017significantly}, and SZ3~\cite{liang2022sz3,zhao2021optimizing}.
Under comparable data distortion, SZ3-HOHS achieves around 1.5$\times$ higher compression ratios than the best-performing baseline, while MRCZ further improves compression ratios by around 5.5$\times$ over SZ3-HOHS. In end-to-end Globus experiments, SZ3-HOHS achieves the lowest total data-movement time among all evaluated methods, while MRCZ remains competitive by minimizing transfer volume despite its additional computational cost.

The contributions of this work are summarized as follows:
\begin{itemize}
\item We highlight high-order hybrid-structured (HOHS) simulation data as an important yet underexplored class of scientific data that presents unique data-management challenges.
\item We design SZ3-HOHS, adapted from SZ3, to provide a high-throughput and storage-efficient option for HOHS data management.
\item We design MRCZ, a model-based solution that maximizes storage reduction.
\item We conduct comprehensive evaluations on large-scale HOHS datasets spanning multiple domains and simulation software packages, measuring storage reduction, end-to-end Globus-based data movement, strong-scaling performance, and reconstruction quality. The results demonstrate that the proposed solutions reduce storage footprint while maintaining competitive end-to-end data-movement performance relative to state-of-the-art error-bounded compressors.
\end{itemize}

The rest of this paper is organized as follows. Section~\ref{sec:related work} reviews error-bounded lossy compression methods for structured and unstructured scientific data, as well as existing compressors specifically designed for HOHS data. Section~\ref{sec:problem formulation} describes the characteristics of HOHS datasets and formulates the data management problem. Section~\ref{sec:design overview} presents the common preprocessing steps for HOHS data and the compression framework that forms the basis of SZ3-HOHS and MRCZ. Sections~\ref{sec:SZ3-HOHS} and \ref{sec:mrcz} present the detailed designs of SZ3-HOHS and MRCZ, respectively. Section~\ref{sec:performance evaluation} evaluates the proposed methods against adapted state-of-the-art compressors on real-world HOHS datasets. Section~\ref{sec:conclusion} concludes the paper.